\documentclass{kaiart}


\usepackage[font=itshape]{quoting}
\quotingsetup{vskip=0pt}

\usepackage{lipsum}
\setlength{\parindent}{0pt}
\setlength{\parskip}{10pt}

\linespread{1.2}

\title{
    A Template for
    Response to Reviewers
}

\newcommand{\cmt}[2]{\textbf{#1}: {\color{gray}{#2}}}

\begin{document}
\maketitle


Dear Editor and Reviewers,


Thanks for your constructive comments and giving us the opportunity
to revise our manuscript (\textbf{Submission-ID}).

Both reviewers acknowledged the novelty and significance of our work, while also offering valuable suggestions for improvement.
%
Most of their comments focused on 
requests for additional experimental results 
and analyses,
as well as clarification of specific 
sections of the manuscript.


We have carefully addressed all comments and 
suggestions and made corresponding revisions 
to improve the manuscript.
%
All revisions in the manuscript are {\color{blue}{highlighted in blue}} for easy identification.


In the rest of the letter, we address each reviewer's comments point-by-point.


\vspace{3em}
Best,

First Name Last Name,

Professor,

School of Computer Science,

University of X


\clearpage
\pagebreak

\section*{Response to Reviewer \#1}
\cmt{R1 Comment1}{
    \lipsum[1]
}

\lipsum[1]

\cmt{R1 Comment2}{
    \lipsum[1]
}

\lipsum[1]

\begin{quoting}
    \lipsum[1]
\end{quoting}

Thanks for the insightful suggestion.

\lipsum[1]


\clearpage\pagebreak
\section*{Response to Reviewer \#2}
\cmt{R1 Comment1}{
    \lipsum[1]
}

\lipsum[1]

\cmt{R1 Comment2}{
    \lipsum[1]
}

\lipsum[1]

\begin{quoting}
    \lipsum[1]
\end{quoting}

Thanks for the insightful suggestion.

\lipsum[1]
~\cite{krizhevsky2012imagenet}

\bibliographystyle{IEEEtran}
\bibliography{eg.bib}

\end{document}